\documentclass[11pt]{article}
\usepackage[margin=1in]{geometry}
\usepackage{amsmath,amssymb,amsthm}
\usepackage[colorlinks=true,linkcolor=blue,urlcolor=blue,citecolor=blue]{hyperref}
\newtheorem{theorem}{Theorem}[section]
\newtheorem{lemma}[theorem]{Lemma}
\newtheorem{corollary}[theorem]{Corollary}
\newtheorem{remark}[theorem]{Remark}
\newcommand{\R}{\mathbb R}
\newcommand{\E}{\mathbb E}
\newcommand{\Prob}{\mathbb P}
\newcommand{\norm}[1]{\left\lVert#1\right\rVert}
\newcommand{\tr}{\operatorname{tr}}
\newcommand{\diag}{\operatorname{diag}}
\newcommand{\sr}{\operatorname{sr}}
\newcommand{\DD}{\mathcal D}
\title{Balanced graph sketches for spectral matrix multiplication}
\author{Working research manuscript}
\date{November 2026}
\begin{document}
\maketitle
\begin{abstract}
We analyze approximate matrix multiplication by averaging random balanced
bucket partitions and correcting their common constant direction exactly.
The bucket incidences are unsigned; the complete implicit sketch includes
a deterministic rank-one mean correction. An exact covariance map retains
a negative semidefinite diagonal term. Combinatorial matrix Bernstein
controls the two bucketed input factors, and a tail-integral truncation
argument gives spectral error $\varepsilon\norm A\norm B$ with
$O(r\varepsilon^{-2}\log^3((r+n+p)/(\delta\varepsilon)))$ rows and
$O(\varepsilon^{-1}\log^3((r+n+p)/(\delta\varepsilon)))$ bucket passes,
where $r$ bounds the two stable ranks. A simpler one-partition bound gives
ordinary $O(N^2)$ total arithmetic, including explicit output, for fixed
stable ranks, accuracy, and confidence. The proof uses established matrix
concentration inequalities. The row/error tradeoffs have stronger signed
sketching precedents, and no asymptotic state-of-the-art or priority claim
is made. The specific sharp mixed-permutation moment analysis motivating
this construction remains a separate route to stronger bounds.
\end{abstract}

\section{Construction, guarantees, and scope}
All matrices are real. The unadorned norm is the operator norm, and
$\sr(A)=\norm A_F^2/\norm A^2$ for nonzero $A$.
Fix $A\in\R^{n\times m}$ and $B\in\R^{m\times p}$ independently of
the sketch randomness. A supplied number
\[
 r\ge\max\{1,\sr(A),\sr(B)\}
\]
will determine the sketch size. If either factor is zero, return zero.
Initially assume $m=sq\ge4$, where $s\ge2$ and $q\ge1$ are integers.
Let $H$ be the $s$-by-$m$ incidence matrix of a uniform balanced partition
of $[m]$ into $s$ labeled buckets, each containing $q$ coordinates.
Such a partition is generated by a uniform permutation followed by a
fixed balanced partition. Put
\begin{equation}\label{eq:W}
 \rho=\frac{q-1}{m-1},\quad \gamma=(1-\rho)^{-1},\quad
 J=\frac{\mathbf1\mathbf1^T}{m},\quad \Pi=I-J,\qquad
 W=\gamma(H^TH-\rho\mathbf1\mathbf1^T).
\end{equation}
Two distinct coordinates collide with probability $\rho$, so
$\E W=I$. In every realization $W\mathbf1=\mathbf1$.

More precisely, with $J_s=\mathbf1_s\mathbf1_s^T/s$ and $\Pi_s=I-J_s$,
\begin{equation}\label{eq:projection}
 P_H=\frac1q H^T\Pi_s H,\qquad
 W=J+\lambda P_H,\qquad \lambda=\frac{m-1}{s-1}.
\end{equation}
Here $P_H$ is an orthogonal projection of rank $s-1$, orthogonal to
constants. Thus $W$ is PSD of rank $s$. Its Gram factor
\begin{equation}\label{eq:singlefactor}
 S=[\sqrt\gamma\,\Pi_s+q^{-1/2}J_s]H,
 \qquad S^TS=W,
\end{equation}
is generally dense because of deterministic mean subtraction. Only $H$
is an unsigned sparse incidence matrix. We do not claim that $S$ has one
nonzero per column or that the complete embedding is nonnegative.

For $d$ independent partitions let
\[
 \overline W=\frac1d\sum_{c=1}^d W_c,\qquad C=A\overline W B.
\]
Then $\E C=AB$. A Gram factor can share the constant direction once,
using $\ell=1+d(s-1)$ rows. Its implicit application uses $d$ bucket
passes and deterministic centering, as detailed in Section~\ref{sec:cost}.

\begin{theorem}[Spectral concentration]\label{thm:main}
Let $0<\varepsilon,\delta<1$, and choose a bucket count satisfying
\begin{equation}\label{eq:s}
 s-1\ge r/\varepsilon.
\end{equation}
There is a universal constant $C_0$ such that, with
\begin{equation}\label{eq:d}
 d=\left\lceil C_0\varepsilon^{-1}
       \log^3\!\left(\frac{C_0(2s+n+p)}{\delta\varepsilon}\right)
      \right\rceil,
\end{equation}
the estimator satisfies
\begin{equation}\label{eq:mainerror}
 \Prob\{\norm{C-AB}>\varepsilon\norm A\norm B\}\le\delta.
\end{equation}
The explicit conservative value $C_0=2^{20}$ is sufficient.
For $s=\lceil1+r/\varepsilon\rceil$ the Gram factor has
\begin{equation}\label{eq:width}
 \ell=O\!\left(r\varepsilon^{-2}
            \log^3\!\frac{r+n+p}{\delta\varepsilon}\right)
\end{equation}
rows and can be applied to the inputs in
$O(dm(n+p)+ds(n+p))$ arithmetic. Explicit expansion adds $O(np\ell)$.
Zero padding and an exact identity fallback remove the divisibility
restriction without increasing these asymptotic bounds.
\end{theorem}

The implicit sufficient inequality~\eqref{eq:implicit} below is less
conservative than~\eqref{eq:d}. We count real arithmetic, comparisons,
positive scalar square roots, and discrete permutation randomness.
Scalar parameter selection has only polynomial logarithmic overhead.
No finite-word numerical stability theorem is asserted.

\begin{corollary}[An ordinary quadratic-output regime]\label{cor:quadratic}
There is also a one-partition version of~\eqref{eq:mainerror} with
\begin{equation}\label{eq:singlechoice}
 d=1,\qquad
 s=\left\lceil1+\frac{3(r^2+2r)}{\varepsilon^2\delta}\right\rceil.
\end{equation}
For $n=p=N$, $m\le N$, and fixed $r,\varepsilon,\delta$, its total
ordinary arithmetic is $O(N^2)$, including every output entry.
The factors may have full algebraic rank; their stable ranks are bounded.
\end{corollary}

The guarantee is absolute error scaled by $\norm A\norm B$.
It is not a $(1+\theta)\sigma_{k+1}(AB)$ low-rank guarantee, or an
error bound scaled by $\norm{AB}$ when factors cancel.
Classical AMM already gives comparable or stronger row/error tradeoffs
\cite{CNW}; the quadratic regime is a construction-specific consequence.

\section{The exact covariance map}
Let $T=W-I$. The identities above give $T\mathbf1=0$, $T_{ii}=0$, and
\[
 \E T^2=c\Pi,\qquad c=\frac{m-s}{s-1}.
\]
Permutation symmetry and the zero row sums determine all covariances of
off-diagonal entries. Write
\begin{equation}\label{eq:coefficients}
 \begin{gathered}
 v=\frac cm,\quad u=-\frac v{m-2},\quad
 w=\frac{2v}{(m-2)(m-3)},\\
 a=v-2u+2w,\quad b=v-2u+w,\quad
 \kappa=2(v-4u+3w).
 \end{gathered}
\end{equation}
The covariance is $v$ for equal unordered coordinate pairs, $u$ for
pairs sharing one vertex, and $w$ for disjoint pairs. Indeed,
$\sum_{j\ne i}\E T_{ij}^2=c(1-1/m)$ gives $v=c/m$.
Taking $\E[T_{ij}\sum_k T_{ik}]=0$ gives $u$, and then the total
row-sum identity gives $w$.

\begin{lemma}[Covariance map]\label{lem:cov}
For every real symmetric $G$ with $G\mathbf1=0$,
\begin{equation}\label{eq:covmap}
 \E[TGT]=aG+b\tr(G)\Pi-\kappa\Pi\diag(G)\Pi.
\end{equation}
Here $\diag(G)$ denotes the diagonal matrix with diagonal entries $G_{ii}$.
\end{lemma}
\begin{proof}
Separate the two summed off-diagonal $T$ entries into the three overlap
classes. Their contribution to diagonal entry $i$ is
\[
 (v-u)\tr G+(-v+2u)G_{ii}.
\]
For $i\ne j$, the corresponding entry is
\[
 (v-2u+2w)G_{ij}+(u-w)[\tr G-2G_{ii}-2G_{jj}].
\]
The row sums of $G$ vanish. Expanding the right side of
\eqref{eq:covmap}, and substituting~\eqref{eq:coefficients}, gives exactly
these two formulas. This proves the full matrix identity.
\end{proof}

Normalize $\norm A=\norm B=1$ for the rest of the proof, and put
$A_0=A\Pi$, $B_0=\Pi B$. The one-color error is
$X=ATB=A_0TB_0$, with $\E X=0$. The lemma gives
\begin{align}\label{eq:leftcov}
 \E[XX^T]
 &=a(A_0B_0)(A_0B_0)^T+b\norm{B_0}_F^2 A_0A_0^T\nonumber\\
 &\hspace{12mm}-\kappa A_0\diag(B_0B_0^T)A_0^T.
\end{align}
The analogous expression holds for $\E[X^TX]$. Since $\kappa\ge0$,
the last term is negative semidefinite. For $m\ge4$,
\begin{equation}\label{eq:abounds}
 b=\frac{c(m-1)}{m(m-3)}\le\frac3{s-1},\qquad
 a=b\left[1+\frac2{(m-1)(m-2)}\right]\le\frac43 b.
\end{equation}
The centered matrices have operator norms at most one and squared
Frobenius norms at most $r$. Consequently
\begin{equation}\label{eq:variance}
 \E[XX^T]\preceq VI_n,\qquad \E[X^TX]\preceq VI_p,
 \qquad V=\frac{3(r+2)}{s-1}.
\end{equation}
Under~\eqref{eq:s}, $V\le9\varepsilon$.

This also proves Corollary~\ref{cor:quadratic}. Take the trace of
\eqref{eq:leftcov}, discard the negative term, and use
$\norm{A_0B_0}_F^2\le r$ to obtain
\[
 \E\norm X_F^2\le ar+br^2
 \le\frac{4r+3r^2}{s-1}\le\frac{3(r^2+2r)}{s-1}.
\]
Markov's inequality and $\norm X\le\norm X_F$ now give
\eqref{eq:mainerror} with~\eqref{eq:singlechoice}.

\section{A tail bound for the bucketed factors}
For a rectangular matrix $Z$, let
\[
 \DD(Z)=\begin{pmatrix}0&Z\\ Z^T&0\end{pmatrix}.
\]
Then $\lambda_{\max}(\DD(Z))=\norm Z$ and
$\DD(Z)^2=\diag(ZZ^T,Z^TZ)$.
The following source theorem is the concentration input.

\begin{lemma}[Combinatorial Bernstein, \cite{Mackey}, Corollary 10.3]
Let $(F_{ij})_{i,j=1}^m$ be deterministic Hermitian $D$-by-$D$ matrices
such that $\sum_{ij}F_{ij}=0$ and $\norm{F_{ij}}\le R$. For a uniform
permutation $\pi$ of $[m]$, put $Z=\sum_iF_{i,\pi(i)}$ and
$\sigma^2=\norm{m^{-1}\sum_{ij}F_{ij}^2}$. Then
\begin{equation}\label{eq:combBernstein}
 \Prob\{\lambda_{\max}(Z)\ge t\}
 \le D\exp\left(-\frac{t^2}{12\sigma^2+4\sqrt2 Rt}\right),\qquad t\ge0.
\end{equation}
\end{lemma}

Let $U\in\R^{m\times r'}$ satisfy $U^T\mathbf1=0$, $\norm U\le1$,
and $\norm U_F^2\le r$. Its columns need not be orthonormal, and $r'$
can exceed $r$. For fixed balanced bucket labels $b(j)$, set
$a_j=e_{b(j)}-\mathbf1_s/s$ and write $u_i^T$ for row $i$ of $U$.
Apply the lemma to $F_{ij}=\DD(a_ju_i^T)$. Their norms are at most one
and their total sum is zero. The permuted sum is the dilation of $HU$.
Moreover,
\begin{equation}\label{eq:combproxy}
 \frac1m\sum_{ij}F_{ij}^2
 =\diag\left(\frac{\norm U_F^2}{s}\Pi_s,
                   (1-1/s)U^TU\right).
\end{equation}
If $s\ge r$, its norm is at most one. Since
$12+64\sqrt2 z\le128z$ for $z\ge1$, putting $t=16z$ gives
\begin{equation}\label{eq:HUtail}
 \Prob\{\norm{HU}\ge16z\}\le(s+r')e^{-2z},\qquad z\ge1.
\end{equation}
The theorem's condition~\eqref{eq:s} implies $s\ge r$.
Apply~\eqref{eq:HUtail} to $U=A_0^T$ and $U=B_0$. A union bound,
$\gamma\le2$, and
\[
 X=\gamma A_0H^THB_0-A_0B_0
\]
show, with $D=2s+n+p$, that
\begin{equation}\label{eq:Xtail}
 \Prob\{\norm X>513z^2\}\le De^{-2z},\qquad z\ge1.
\end{equation}
This step has no ambient $m/s$ norm factor and assumes no entrywise
incoherence. It uses the Frobenius bounds of the input factors directly.

\section{Truncation and spectral concentration}
Fix a positive integer $d$ and define
\begin{equation}\label{eq:truncparams}
 h=\log\frac{8192Dd}{\delta\varepsilon},\qquad \mathcal R=513h^2.
\end{equation}
Here $h\ge1$. For independent colors, let
$Y_c=X_c\mathbf1_{\{\norm{X_c}\le\mathcal R\}}$ and $B_* =\E Y_c$.
Truncation is only an analysis device. The algorithm returns the
untruncated average.

\begin{lemma}[Bias and covariance after truncation]\label{lem:trunc}
The bias obeys
\begin{equation}\label{eq:bias}
 \norm{B_*}\le\frac{\delta\varepsilon}{4d},
\end{equation}
and the probability that any of the $d$ colors is truncated is at most
$\delta/4$. The independent centered dilations
$Z_c=\DD(Y_c-B_*)$ have norm at most $2\mathcal R$ and
$\E Z_c^2\preceq VI_{n+p}$.
\end{lemma}
\begin{proof}
Since $\E X=0$, the tail-integral identity and~\eqref{eq:Xtail} give
\begin{align*}
 \norm{B_*}
 &\le\E[\norm X\mathbf1_{\{\norm X>\mathcal R\}}]\\
 &\le513De^{-2h}(h^2+h+1/2)
 \le1283Dh^2e^{-2h}\le1283De^{-h}
 \le\delta\varepsilon/(4d).
\end{align*}
For the penultimate inequality use $h^2e^{-h}\le1$ for $h\ge1$.
In detail, substituting $t=513z^2$ bounds the integrated tail by
\[
 \int_{\mathcal R}^{\infty}\Prob\{\norm X>t\}\,dt
 \le1026D\int_h^{\infty}ze^{-2z}\,dz
 =513De^{-2h}(h+1/2).
\]
The other $h^2$ term is $\mathcal R\Prob\{\norm X>\mathcal R\}$.
A union bound yields truncation probability at most
$dDe^{-2h}\le\delta/4$. Finally, write
$\widehat X=\DD(X)$ and $\widehat Y=\DD(Y)$. Then
\[
 \E[(\widehat Y-\E\widehat Y)^2]
 =\E\widehat Y^2-(\E\widehat Y)^2
 \preceq\E\widehat Y^2\preceq\E\widehat X^2\preceq VI.
\]
The norm bound follows from $\norm Y\le\mathcal R$ and
$\norm{B_*}\le\varepsilon/4\le\mathcal R$.
\end{proof}

We also invoke the independent matrix Bernstein inequality in
\cite[Corollary 5.2]{Mackey}: independent mean-zero Hermitian matrices
with norms at most $R$ and sum of second moments of norm at most
$\sigma^2$ satisfy
\[
 \Prob\{\lambda_{\max}(\textstyle\sum_cZ_c)\ge t\}
 \le(n+p)\exp\left(-\frac{t^2}{3\sigma^2+2Rt}\right).
\]
Apply it with $R=2\mathcal R$, $\sigma^2\le dV$, and
$t=d\varepsilon/2$. Because the matrices remain dilations, this gives
\begin{equation}\label{eq:trimtail}
 \Prob\left\{\norm{\frac1d\sum_c(Y_c-B_*)}>\varepsilon/2\right\}
 \le(n+p)\exp\left(-\frac{d\varepsilon^2}
                              {12V+8\mathcal R\varepsilon}\right).
\end{equation}
It therefore suffices to choose
\begin{equation}\label{eq:implicit}
 d\ge\left[\frac{12V}{\varepsilon^2}
              +\frac{4104h^2}{\varepsilon}\right]
               \log\frac{4(n+p)}\delta.
\end{equation}
Off the two bad events, the original average has norm at most
$\varepsilon/2+\norm{B_*}<\varepsilon$. Their total probability is
at most $\delta/2$, so~\eqref{eq:mainerror} follows. Restoring the two
input norms supplies the claimed error scale.

To verify the explicit choice, put
$L=\log(C_0D/(\delta\varepsilon))$, $C_0=2^{20}$, and
$d=\lceil C_0\varepsilon^{-1}L^3\rceil$. Then $L>13.8$ and
$d\le2C_0\varepsilon^{-1}L^3$, giving
\[
 \log d<3L,\qquad h<4L,\qquad
 \log(4(n+p)/\delta)\le L.
\]
Together with $V\le9\varepsilon$, the right side of
\eqref{eq:implicit} is at most
\[
 \varepsilon^{-1}(108+65664L^2)L
 <C_0\varepsilon^{-1}L^3\le d.
\]
This proves Theorem~\ref{thm:main}. The large displayed constant favors
an elementary check; the implicit test~\eqref{eq:implicit} is preferable
when selecting concrete parameters.

\section{Complete factors, padding, and arithmetic}\label{sec:cost}
The one-partition construction in Corollary~\ref{cor:quadratic} can be
implemented by the following pseudocode. The variables $a_*$ and $b_*$
here are scalar factor coefficients, separate from the covariance
coefficients $a,b$ above. Bucket means are taken over all $s$ buckets.
\begin{quote}\small
\begin{verbatim}
Input A, B, a stable-rank bound r, epsilon, delta
s = ceil(1 + 3*(r*r + 2*r)/(epsilon*epsilon*delta))
If s >= m, return the exact product A B
M = s*ceil(m/s)
Pad A and B with zero columns and rows to inner dimension M
Draw a uniform permutation and split it into s equal buckets
For each bucket j, sum its A columns into U[:,j]
For each bucket j, sum its B rows into V[j,:]
q = M/s, rho = (q-1)/(M-1)
a_star = 1/sqrt(1-rho), b_star = 1/sqrt(q)
X[:,j] = a_star*U[:,j] + (b_star-a_star)*mean_columns(U)
Y[j,:] = a_star*V[j,:] + (b_star-a_star)*mean_rows(V)
Return C = X Y, forming every entry by ordinary multiplication
\end{verbatim}
\end{quote}
The repeated mean in each displayed assignment is computed once.
These formulas apply $F=\sqrt\gamma\Pi_s+q^{-1/2}J_s$ to the bucketed
inputs: $X=(AH^T)F$, $Y=F(HB)$. Thus their product is exactly $AWB$,
without materializing a dense sketch. Permutation generation costs
$O(M)$ and zero padding can be represented implicitly.

For multiple colors, the projection representation~\eqref{eq:projection} gives
\[
 \overline W=J+\frac\lambda d\sum_cP_{H_c}.
\]
Choose one fixed Householder reflection in bucket coordinates that sends
$\mathbf1_s/\sqrt s$ to $e_1$. Apply it to $\Pi_sH_c$ and delete its
identically zero first row, obtaining $\widetilde H_c$. Then
$\widetilde H_c\widetilde H_c^T=qI_{s-1}$ and
$P_{H_c}=q^{-1}\widetilde H_c^T\widetilde H_c$.
An explicit shared-channel factor is
\[
 S_*=
 \begin{bmatrix}
  \mathbf1^T/\sqrt m\\
  \sqrt{\gamma/d}\,\widetilde H_1\\
  \vdots\\
  \sqrt{\gamma/d}\,\widetilde H_d
 \end{bmatrix},\qquad S_*^TS_*=\overline W.
\]
The returned factors are $AS_*^T$ and $S_*B$.
Computing the bucket sums costs $O(dm(n+p))$. Centering and the fixed
Householder transforms require $O(ds(n+p))$, using a dot product and
rank-one update per vector. There is no dense row orthogonalization.
All factor entries use $O(\ell(n+p))$ storage, and naive expansion costs
$O(np\ell)$ arithmetic and $O(np)$ output storage.

If $s<m$ does not divide $m$, pad by zeros to
$M=s\lceil m/s\rceil<m+s\le2m$ and apply the theorem at $M$.
The operator norms, Frobenius norms, stable ranks, and exact product
are unchanged. If $s\ge m$, return the exact factors $A,B$ instead,
with width $m\le s$. Small ambient dimensions can also use this fallback.
One may additionally choose the identity whenever its width is smaller
than the proposed $1+d(s-1)$; this strengthens the bound without altering
the proof. The costs in Corollary~\ref{cor:quadratic} now follow directly.

All random incidence operations are additions of input entries.
The complete $S_*$ includes deterministic negative entries from mean
subtraction. Thus the proved property is fast application of an implicit
sparse-plus-low-rank map, not a literal per-column sparsity guarantee.
Integer parameter choices can use comparisons and doubling, with
logarithmic scalar approximations from standard convergent series.
Their overhead is absorbed by the stated construction work.

\section{Prior work and remaining research boundary}
Cohen--Nelson--Woodruff~\cite{CNW} establish stable-rank spectral AMM
with optimal $O(r/\varepsilon^2)$ row dependence. The present logarithmic
row bound is weaker. For OSEs, the fully independent signed SparseStack
construction is reported with optimal row count
$O((r+\log(1/\delta))/\varepsilon^2)$ and
$O(\log(r/\delta)/\varepsilon)$ nonzeros per column~\cite{SparseStack}.
Our application-time bound concerns the incidence passes and deterministic
mean correction, and is not an improvement to those sparsity results.

The covariance calculation exploits regular bucket sizes and permutation
row sums. Its full higher-order trace expansion consists of the colored
centered permutation moments studied in the companion
\texttt{spectral-gap} project. That project supplies exact mixed-cell
envelopes, but no use of their sharper high-order constants is claimed
in this Bernstein proof. The source theorem~\eqref{eq:combBernstein}
uses the entire dependent permutation law directly.

The contribution of this derivation is a complete concentration argument
for the particular balanced-partition estimator, its stable-rank variance
map, full factor and output costs, and explicit parameter accounting.
No first-discovery claim is made for this synthesis. It proves neither
general quadratic matrix multiplication nor the stronger arbitrary-factor
relative low-rank target. Those goals need further mathematical advances.

\begin{thebibliography}{9}
\bibitem{Mackey} L. Mackey, M. I. Jordan, R. Y. Chen, B. Farrell, J. A. Tropp.
Matrix concentration inequalities via the method of exchangeable pairs.
Annals of Probability 42 (2014), 906--945.
Corollaries 5.2 and 10.3.
\url{https://lmackey.github.io/papers/matstein-aop14.pdf}.
\bibitem{CNW} M. B. Cohen, J. Nelson, D. P. Woodruff.
Optimal Approximate Matrix Product in Terms of Stable Rank.
ICALP 2016, LIPIcs 55, 11:1--11:14.
\url{https://doi.org/10.4230/LIPIcs.ICALP.2016.11}.
\bibitem{SparseStack} D. Heidary.
SparseStack Is an Optimal Oblivious Subspace Embedding.
Preprint, 2026, version 2.
\url{https://arxiv.org/abs/2609.02978}.
\end{thebibliography}
\end{document}
